
(iii)$\To$(ii).
\marginpar{65}
Sea~$\goth{n}$ el ideal maximal de~$C$ y, para todo entero $q>0$, sea
$$
C_q=C/\goth{n}^q,\qquad J_q=(J+\goth{n}^q)/\goth{n}^q.
$$
Así, $C_q/J_q$ se identifica con un álgebra cociente de~$C/J$; por otra parte, el homomorfismo compuesto $u_q\colon B\to C/J\to C_q/J_q$ es continuo de~$B$ en el anillo discreto~$C_q/J_q$, y~$J_q$ es un ideal nilpotente de~$C_q$. Se construyen sucesivamente $A$-homomorfismos
$$
v_q\colon B\to C_q
$$
tales que (a)~$v_q$ levanta~$u_q$ y (b)~$v_q$ levanta~$v_{q-1}$. La posibilidad de la inducción se comprueba fácilmente, pues, como
$$
u_q\colon B\to C/(J+\goth{n}^q)\qquad\text{y}\qquad
v_{q-1}\colon B\to C/\goth{n}^{q-1}
$$
definen el mismo homomorfismo
$$
B\to C/((J+\goth{n}^q)+\goth{n}^{q-1})
=C/(J+\goth{n}^{q-1})=C_{q-1}/J_{q-1},
$$
a saber, $u_{q-1}$, determinan un homomorfismo
$$
B\to C/J_q'\qquad\text{donde}\qquad
J_q'=(J+\goth{n}^q)\cap\goth{n}^{q-1}\supset\goth{n}^q
$$
(del cual ambos proceden por reducción). Se reduce, pues, el problema a levantar un homomorfismo $B\to C/J_q'$ de~$B$ en un cociente de~$C_q$ por un ideal $J_q'/\goth{n}^q$ contenido en~$J_q$ y, por tanto, nilpotente; ello es posible por la hipótesis~(iii).

Hecho esto, los~$v_q$ definen un homomorfismo de~$B$ en el límite proyectivo~$C$ de los~$C_q$. Como~$J$ es cerrado, $J$ es el límite proyectivo de los~$J_q$; por consiguiente, $v$ levanta~$u$, como se quería demostrar.

(iv)$\To$(i). Se comprueba primero de inmediato que, si se verifica~(iv), sigue verificándose para los componentes locales de $B\otimes_A A'$ sobre~$A'$, cuando~$A'$ es una $A$-álgebra local finita. Tomando~$A'$ libre sobre~$A$ y tal que las extensiones residuales de~$B'$ sobre~$A'$ sean triviales, se reduce el problema, gracias a~\Ref{III.1.4}~(ii), al caso en que la extensión residual de~$B$ sobre~$A$ es trivial. Demostraremos entonces el resultado algo más preciso siguiente.

\begin{corollaire}
\label{III.2.2}
Bajo las condiciones de~\Ref{III.2.1}, supongamos además que la extensión residual de~$B$ sobre~$A$ es trivial. Entonces las condiciones equivalentes
\marginpar{66}
de~\Ref{III.2.1} son también equivalentes a las dos condiciones siguientes (suponiendo en \textup{(v)} que~$A$ y~$B$ son completos):
\begin{enumerate}
\item[(iv~ter)] Como \textup{(iv)}, pero restringiendo el anillo local artiniano~$C$, finito sobre~$A$, a tener una extensión residual trivial (y, además, si se desea, suponiendo que el ideal~$J$ es de cuadrado nulo).
\item[(v)] Existe un $A$-homomorfismo local (donde $n=\dim \goth{n}/(\goth{n}^2+\goth{m}B)$)
$$
u\colon B\to B_1=A[[t_1,\dots,t_n]]
$$
que induce un \emph{isomorfismo}
$$
\goth{n}/(\goth{n}^2+\goth{m}B)\isomto
\goth{n}_1/(\goth{n}_1^2+\goth{m}B_1),
$$
donde~$\goth{n}$ (respectivamente $\goth{n}_1$) es el ideal maximal de~$B$ (respectivamente de~$B_1$), y~$\goth{m}$ el de~$A$.
\end{enumerate}
\end{corollaire}

\subsubsection*{Demostración}
Como (iv~bis) implica evidentemente (iv~ter) ---dejando de lado la broma del ideal de cuadrado nulo---, bastará demostrar (iv~ter)$\To$(v)$\To$(i).

(iv~ter)$\To$(v). Elijamos una base $a_1,\dots,a_n$ de $\goth{n}/(\goth{n}^2+\goth{m}B)$, lo que define un homomorfismo local de $A$-álgebras
$$
B\to B_1/(\goth{n}_1^2+\goth{m}B_1)
=k[t_1,\dots,t_n]/(t_1,\dots,t_n)^2,
$$
que, en virtud de~(iv~ter), puede levantarse sucesivamente a homomorfismos de $A$-álgebras de~$B$ en $B_1/\goth{n}_1^2$, $B_1/\goth{n}_1^3$, etc.; pasando al límite proyectivo se obtiene el homomorfismo $B\to B_1$ con la propiedad deseada.

(v)$\To$(i). Como, en el diagrama conmutativo
$$
\ifthenelse{\boolean{orig}}
{\xymatrix{ \goth{m}/\goth{m}^2 \ar[r] & \goth{n}/\goth{n}^2 \ar[r]
\ar[d] & \goth{n}/(\goth{n}^2+\goth{m}B) \ar[r] \ar[d] & 0 \\
\goth{m}/\goth{m}^2 \ar[r] & \goth{n}_1/\goth{n}_1^2 \ar[r] &
\goth{n}_1/(\goth{n}_1^2+\goth{m}B_1) \ar[r] & 0}}
{\xymatrix{ \goth{m}/\goth{m}^2 \ar[r] \ar[d] & \goth{n}/\goth{n}^2
\ar[r] \ar[d] & \goth{n}/(\goth{n}^2+\goth{m}B) \ar[r] \ar[d] & 0 \\
\goth{m}/\goth{m}^2 \ar[r] & \goth{n}_1/\goth{n}_1^2 \ar[r] &
\goth{n}_1/(\goth{n}_1^2+\goth{m}B_1) \ar[r] & 0}}
$$
las dos filas son exactas y las flechas verticales extremas son sobreyectivas, la flecha intermedia es sobreyectiva; de ello se sigue (al ser~$B$ completo) que $B\to B_1$ es sobreyectivo. Sean~$x_i$ ($1\leq i\leq n$) elementos de~$B$ que levantan los~$t_i$. Definen un homomorfismo de $A$-álgebras $B_1\to B$, que será sobreyectivo por la misma razón que~$u$ y cuya composición con~$u$ es la identidad por construcción. Por tanto, $B_1\to B$ es también inyectivo y, en consecuencia, un isomorfismo. Se obtiene así el corolario siguiente.

\begin{corollaire}
\label{III.2.3}
Bajo
\marginpar{67}
las condiciones de~\Ref{III.2.2}~\textup{(v)}, $u$ es necesariamente un isomorfismo.
\end{corollaire}

Esto concluye la demostración de que~$B$ es formalmente lisa sobre~$A$. Al mismo tiempo hemos recuperado~\Ref{III.1.5} (aunque no hay gran mérito en ello).
